\documentclass[10pt,a4paper]{amsart}
\usepackage[margin=19mm]{geometry}
\usepackage{amsmath,amssymb,mathtools}
\usepackage[scr=rsfs,cal=euler]{mathalfa}
\usepackage{xfrac}
\usepackage[unicode,hidelinks,bookmarks=false]{hyperref}
\newcommand{\ie}{i.e.\ }
\newcommand{\cf}{cf.\ }
\newcommand{\resp}{respectively\ }
\newcommand{\Subsection}{\subsection}
\providecommand{\nobreakdash}{\nobreak\hbox{-}}
\setlength{\emergencystretch}{2em}
\multlinegap0pt
\usepackage{amssymb}
\usepackage{mathtools}
\usepackage[shortlabels]{enumitem}
\newtheorem{lemma}{Lemma}
\newcommand{\cA}{\mathcal A}
\newcommand{\cB}{\mathcal B}
\title{A nine-line counterexample to a conjecture on the minimal degree of Jacobian relations}
\author{Alexandru Dimca, Piotr Pokora}
\date{Source: arXiv:2607.01985v2 (CC BY 4.0)}
\begin{document}
\maketitle

\noindent\emph{Source and licence.} A nine-line counterexample to a conjecture on the minimal degree of Jacobian relations. Version arXiv:2607.01985v2 (CC BY 4.0), distributed by arXiv under the Creative Commons Attribution 4.0 International licence, http://creativecommons.org/licenses/by/4.0/, as recorded in the arXiv licence metadata for this exact version and shown on the abstract page https://arxiv.org/abs/2607.01985v2. Full licence notice: \texttt{licenses/arxiv-2607.01985-CC-BY-4.0.txt}.\par\smallskip

\noindent\textit{Editorial scope.} Counted unit: the article's lemma lem:lattice (source0.tex lines 231--246). The article's complete original proof of it is delivered with it (source0.tex lines 248--252, one \texttt{proof} environment of 5 lines). No mathematics is written, repaired or re-derived: the statement and the proof are the article's own lines, in the article's own notation, and no retained line is substituted or re-flowed.  No further source-local context is needed: every cross-reference of every retained line resolves inside the two delivered spans.  Nothing was omitted from any retained span, so the package carries no omission record.  No retained line contains a citation, so the wrapper carries no bibliography and no bibliography entry.  No retained line contains a figure environment, an image-inclusion command or drawing code, so no image is delivered and the package declares no asset dependency.  Editorial formatting only, in addition: an AMS \texttt{amsart} preamble that restates the article's own declarations and macros used by the retained lines, namely the environments \texttt{lemma} on separate counters, together with the standard packages those lines need.  The retained environments print in this document as Lemma 1 in order of appearance, whereas the article prints the same items with the numbering of its own full document.  The complete native source of the article as distributed in the arXiv e-print for this exact version is bundled unchanged as \texttt{sources/204285/source0.tex}.
\begin{lemma}\label{lem:lattice}
With the labelling above, the non-double intersection points of both arrangements are exactly
\begin{equation}\label{eq:incidence-list}
\begin{gathered}
\{1,2,3\},\quad
\{1,4,5\},\quad
\{1,6,7\},\quad
\{2,4,6\},\\
\{3,5,7\},\quad
\{3,6,8\},\quad
\{4,7,9\},\quad
\{2,5,8,9\}.
\end{gathered}
\end{equation}
Consequently, \(\cA\) and \(\cB\) have isomorphic intersection lattices.
\end{lemma}

\begin{proof}
Represent a line \(a x+b y+c z=0\) by the vector \((a,b,c)\).  Three labelled lines \(i,j,k\) are concurrent if and only if the determinant of the corresponding \(3\times3\) matrix is zero.  The same determinant calculation for the two lists of line vectors gives exactly the subsets in \eqref{eq:incidence-list}.  The unique subset of cardinality four is \(\{2,5,8,9\}\), and every other dependent triple is one of the seven triples displayed above.  All other pairs meet in ordinary double points.

Thus the map \(L_i\mapsto M_i\) identifies the rank-two flats of the two arrangements and induces an isomorphism of intersection lattices.
\end{proof}

\end{document}
